\subsection*{SRV Trace 8: Operatio De-Parameterization Correspondence}
\label{subsec:appB_srv_trace8}

\begin{quote}
\emph{``The Operatio is not a poem about the equations.
It is the equations with their parameters removed.''}
\end{quote}

\noindent
This trace demonstrates that the symbolic prelude (the \hyperref[sec:bk1_operatio]{Operatio})
encodes the structural skeleton of PS dynamics. We proceed in two directions:
first stripping the core equations to reveal the Operatio's symbols, then
specializing the Operatio back to recover those equations.
This is a \emph{reflexive validation} --- a correspondence demonstration,
not a retroactive derivation of the axiomatic content of Books~I--IX from
the Operatio alone.

\subsubsection*{8.1 The Core Equations}

The dynamical framework of PS rests on the following constructions from Book~I
and Book~II:

\begin{enumerate}
\item \textbf{Fokker--Planck equation}
(Thm.~\ref{theorem:bk1_fundamental_relation_fokker_plank_equation}):
\[
\frac{\partial \rho}{\partial s}
= -\nabla \cdot (\rho\, D) + \sigma^2 \nabla^2 \rho
\]
governing the evolution of symbolic probability density $\rho$ under drift $D$
and diffusion $\sigma^2$ on the symbolic manifold $M$.

\item \textbf{Symbolic Hamiltonian}
(Def.~\ref{definition:bk2_symbolic_hamiltonian}):
\[
H(x) = \frac{\kappa}{\|D(x)\|_g + \epsilon} + \lambda \cdot \mathrm{tr}(\mathcal{L}_x)
\]
encoding the interplay of drift magnitude and linearized reflection.

\item \textbf{Drift--reflection interplay:} the drift field $D$
(Def.~\ref{definition:bk1_drift_field}), the reflection operator $R$
(Def.~\ref{definition:bk1_reflection_operator}), and the symbolic flow
$\Phi_s$ (Def.~\ref{definition:bk1_symbolic_flow}) on the manifold $M$
(Def.~\ref{definition:bk1_symbolic_manifold}).
\end{enumerate}

\subsubsection*{8.2 De-Parameterization: Equations $\to$ Operatio}

We now factor out all parametric content --- the manifold $M$, metric $g$,
density $\rho$, specific drift field $D$, reflection map $R$, and constants
$\kappa, \lambda, \epsilon, \sigma$ --- retaining only the structural
relations between operations. The correspondence is as follows:

\begin{center}
\renewcommand{\arraystretch}{1.4}
\begin{tabular}{@{}lll@{}}
\textbf{Operatio symbol} & \textbf{Structural skeleton} & \textbf{Parametric source} \\
\hline
$\mathcal{U}_{\emptyset} \;\vdash\; \partial$
& Undifferentiated ground yields differentiation
& FP: $\rho$ evolves under $\nabla \cdot (\rho D)$ \\
$\partial \;\nvdash\; \mathcal{U}_{\emptyset}$
& Differentiation cannot recover the ground
& H-theorem: $dF/ds \leq 0$ (irreversibility) \\
$\mathbb{S} := \mathcal{R}(\partial)$
& Stable structure via reflection on drift
& Equilibrium: $\rho_{\mathrm{eq}} \propto e^{-\beta H}$ \\
$\mathbb{S} \neq \emptyset$
& Viability: stable structure is non-trivial
& Viability domain
(Def.~\ref{definition:bk5_viability_domain}) \\
$\mathcal{O} := \mathcal{E}(\mathbb{S})$
& Observer emerges as convergent identity
& Convergent identity
(Thm.~\ref{theorem:bk7_reflective_convergence_to_stable_identity}) \\
\end{tabular}
\end{center}

\noindent
Each line of the Operatio corresponds to a structural invariant that survives
the removal of all parametric detail. The Operatio is the
\emph{pre-parametric form} of PS dynamics.

\subsubsection*{8.3 Re-Parameterization: Operatio $\to$ Equations}

Conversely, observer specialization recovers the full equations:

\begin{enumerate}
\item \textbf{$\mathcal{U}_{\emptyset} \;\vdash\; \partial$}
(Bounded observer on void $\to$ smooth manifold):
A bounded observer
(Def.~\ref{definition:bk1_bounded_observer}) encountering undifferentiated
ground induces, through successive resolution stages $P_\lambda$, the
colimit manifold $M$ (Thm.~\ref{theorem:bk1_manifold_emergence}).
Differentiation $\partial$ becomes the smooth structure on $M$.

\item \textbf{$\partial \;\nvdash\; \mathcal{U}_{\emptyset}$}
(Irreversibility):
Once differentiation has occurred, the H-theorem
(Thm.~\ref{theorem:bk1_h_theorem_for_symbolic_evolution}) guarantees
monotone free-energy decrease: the undifferentiated state is not recoverable.

\item \textbf{$\mathbb{S} := \mathcal{R}(\partial)$}
(Reflection stabilizes drift):
The drift field $D$ on $M$ is regulated by the reflection operator $R: M \to M$.
Their balance, under smoothness and boundedness, determines $H$
(Def.~\ref{definition:bk2_symbolic_hamiltonian}). The equilibrium density
$\rho_{\mathrm{eq}}$ satisfying $\partial \rho / \partial s = 0$ in the
Fokker--Planck equation yields the Gibbs measure
(Thm.~\ref{theorem:bk2_equilibrium_distribution}).

\item \textbf{$\mathbb{S} \neq \emptyset$}
(Viability):
The viability domain
(Def.~\ref{definition:bk5_viability_domain}) is non-empty: the system
sustains structure under drift--reflection balance. Probability conservation
plus drift plus diffusion recovers the full Fokker--Planck equation.

\item \textbf{$\mathcal{O} := \mathcal{E}(\mathbb{S})$}
(Convergent identity):
The reflective convergence theorem
(Thm.~\ref{theorem:bk7_reflective_convergence_to_stable_identity})
establishes that stable identity $\mathcal{O}$ emerges as a fixed point
of the SRMF cycle acting on viable symbolic structure.
\end{enumerate}

\subsubsection*{8.4 Framing}

This correspondence validates that the Operatio encodes the structural skeleton
of PS dynamics. The axiomatic method of Books~I--IX remains the rigorous
development; the Operatio is neither a substitute for it nor a derivation of it.
This trace demonstrates reflexive consistency: the symbolic prelude and the
formal development describe the same architecture at different levels of
parametric resolution.

\begin{flushright}
\emph{The operation precedes the parameter.\\
The parameter instantiates the operation.\\
Neither is prior; both are necessary.}
\end{flushright}
